\sezione{2}{Entità}{slide 38–48}

\begin{inbreve}
Un'\hl{entità} rappresenta una classe di oggetti della realtà di interesse (persone, cose, fatti) con proprietà comuni ed esistenza autonoma. Nello schema si disegna come un \key{rettangolo} con il nome dentro. Lo schema rappresenta le classi, non i singoli oggetti.
\end{inbreve}

\subsection{Definizione}

\begin{tecnica}[Entità]
Classe di oggetti (fatti, persone, cose) della realtà di interesse con \hl{proprietà comuni} e con \hl{esistenza «autonoma»}, nel senso che i singoli oggetti sono in qualche modo riconoscibili.

Esempi: studente, docente, corso, aula, dipartimento.
\end{tecnica}

\par\noindent\begin{minipage}{\linewidth}
\begin{schema}[x=1cm]
  \foreach \x/\n in {0/Studente,3/Docente,6/Corso,9/Aula,12/Dipartimento}{\node[ent,minimum width=2.4cm] at (\x,0) {\n};}
\end{schema}
\didascalia{Rappresentazione grafica delle entità: un rettangolo con il nome (slide 44).}
\end{minipage}\par

\subsection{Schema e istanza}

\begin{intuizione}[Lo stampo e i biscotti]
L'entità è lo \hl{stampo}: dice che forma hanno tutti gli studenti. Le \key{occorrenze} sono i biscotti: Mario Rossi, Piero Neri e così via. Lo schema concettuale disegna gli stampi, non i biscotti.
\end{intuizione}

\begin{tecnica}[Entità e occorrenza]
\begin{itemize}
  \item \hl{Entità}: classe di oggetti, persone, … «omogenei».
  \item \hl{Occorrenza} di entità: un elemento della classe, cioè \key{l'oggetto o la persona stessa}, non i dati che la descrivono.
  \item Nello schema concettuale si rappresentano le entità, non le singole occorrenze: è un'\key{astrazione}.
\end{itemize}
\end{tecnica}

\begin{attenzione}[Occorrenza non è riga]
Un'occorrenza è la persona Mario Rossi, non la riga di una tabella che contiene i suoi dati. Nel modello relazionale ogni riga «rappresenta» un oggetto; nel modello E-R l'occorrenza \emph{è} l'oggetto. La distinzione aiuta a capire perché due occorrenze diverse possono avere valori uguali per alcuni attributi (due studenti omonimi).
\end{attenzione}

\subsection{Quali entità? L'esempio della biblioteca}

Le specifiche dell'esercizio d'esame del 13/09/2002 (lo risolveremo per intero nel capitolo 9) parlano di prestiti di «libri». Ma che cosa è un libro? Proviamo a chiederci quali sarebbero le \hl{occorrenze} di un'entità \emph{Libro}:
\begin{itemize}
  \item le \key{opere} («I Miserabili» di Victor Hugo)?
  \item le \key{edizioni} (Mondadori, collana Oscar, 1975)?
  \item i \key{volumi} (il terzo volume di quell'edizione)?
  \item le \key{copie} (la seconda copia di quel volume, che sta sullo scaffale)?
\end{itemize}

\par\noindent\begin{minipage}{\linewidth}
\begin{schema}[x=1cm]
  \node[ent] at (0,0) {Opera};
  \node[ent] at (3.2,0) {Edizione};
  \node[ent] at (6.4,0) {Volume};
  \node[ent] at (9.6,0) {Copia};
\end{schema}
\didascalia{Probabilmente quattro entità diverse (slide 47): ognuna ha occorrenze diverse.}
\end{minipage}\par

\begin{esempio}[Contare le occorrenze]
Una biblioteca possiede 2 copie del terzo volume de «I Miserabili» (edizione Mondadori Oscar 1975, in 4 volumi) e 1 copia di ciascuno degli altri tre volumi. Allora ci sono: 1 opera, 1 edizione, 4 volumi e $2+3\cdot1=5$ copie. Un'unica entità \emph{Libro} non saprebbe dire «quanti libri» ci sono: 1, 4 o 5? Il fatto che le risposte siano diverse è il segnale che i concetti sono diversi.
\end{esempio}

\subsection{Convenzioni sui nomi}

\begin{consiglio}[Nomi delle entità]
Ogni entità ha un nome che la identifica \hl{univocamente} nello schema. Usa nomi \key{espressivi} e convenzioni uniformi: nomi al \key{singolare} (\emph{Studente}, non \emph{Studenti}), perché il nome descrive la classe, cioè com'è fatta ciascuna occorrenza.
\end{consiglio}

\begin{esercizi}
\begin{passi}
  \item In una realtà ospedaliera, quali entità useresti per «letto», «reparto», «ricovero»? Ricovero è un'entità o un legame tra altre entità?
  \item Perché \emph{Studenti} non è un buon nome di entità?
\end{passi}
\end{esercizi}

\begin{appunti}{Riepilogo del capitolo}
\pag{39} Entità = classe di oggetti con proprietà comuni ed esistenza autonoma.
\pag{43} Entità (classe) $\neq$ occorrenza (l'oggetto, non i dati).
\pag{47} Un «libro» può essere opera, edizione, volume o copia: concetti diversi.
\pag{48} Nomi espressivi, al singolare, unici nello schema.
\end{appunti}
